%%%%%%%%%%%%%%%%%%%%%%%%%%%%%%%%%%%%%%%%%%%%%%%%%%%%%%%%%%%%%%%%%%%%%%%%%%%%%%%%
%%% THEORY & MANUAL FOR 'ALKCALC'                                            %%%
%%%                                                                          %%%
%%% AUTHOR OF THIS FILE: SIMON EUCHNER                                       %%%
%%%%%%%%%%%%%%%%%%%%%%%%%%%%%%%%%%%%%%%%%%%%%%%%%%%%%%%%%%%%%%%%%%%%%%%%%%%%%%%%

\documentclass[fontsize=12pt,
               twoside=false,
               BCOR=0mm,
               DIV=14,
               overfullrule=true, % SET 'TRUE' FOR OVERFULL BOX INDICATORS
               %chapterprefix=false, % NOT NEEDED FOR SCRARTCL
               toc=bibliography,
               headinclude=true,
               footinclude=false]{scrartcl}

%%% FONT
\usepackage{libertine}
\usepackage{libertinust1math}
\usepackage{dutchcal}
\KOMAoption{DIV}{last}
\setkomafont{captionlabel}{\usekomafont{descriptionlabel}}
\usepackage{upgreek,textcomp}
\AfterTOCHead[toc]{\sffamily}

%%% CAPTIONS
\setcapindent{0em}
\setcapwidth[c]{\textwidth}
\setkomafont{caption}{\sffamily}

%%% LANGUAGE
\usepackage[UKenglish]{babel}

%%% HEADER/FONT
\usepackage[headsepline=1pt:head,
            footsepline=1pt:foot]{scrlayer-scrpage}

%%% FOOTNOTES
\setfootnoterule[.5pt]{.5\textwidth}
\deffootnote{1em}{1em}{
    \makebox[1em][l]{\thefootnotemark}
}
\setkomafont{footnote}{\sffamily}
\setlength{\skip\footins}{1.5em}
\setlength{\footnotesep}{1em}

%%% SECTION AND SUBSECTION SETTINGS
\renewcommand{\thesection}{\Roman{section}}
\renewcommand{\thesubsection}{\roman{subsection}}

%%% GRAPHICS
\usepackage[pdftex]{graphicx}
\usepackage{tikz,color,xcolor}
\graphicspath{{./}}

%%% TABLES
\usepackage{booktabs}

%%% ALIGN TEXT
\usepackage{enumerate}

%%% DOCUMENTATION
\usepackage{sdoc}

%%% EQUATION TYPESETTING
\usepackage{amsmath,amssymb,amsthm,mathtools}
\usepackage{physics,bm,slashed,cancel,tcolorbox}
\numberwithin{equation}{section}

%%% HYPERLINKS
\usepackage[unicode=true,
            colorlinks=false, % SET TO 'FALSE' FOR COLOURED REF'S, ELSE 'TRUE'
            citebordercolor=blue,
            citecolor=black,
            urlbordercolor=blue,
            urlcolor=black,
            linkbordercolor=orange,
            linkcolor=black,
            pdfborder={0 0 1},
            pdfborderstyle={/S/U/W 1}]{hyperref}

%%% BIBLIOGRAPHY
\usepackage{csquotes,bib}
\addbibresource{./references.bib}

%%% ACRONYMS
\usepackage{acro}
\DeclareAcronym{fem}{
    short=FEM,
    long=finite element method,
    short-plural=s,
    long-plural=s,
}
\DeclareAcronym{fdm}{
    short=FDM,
    long=finite difference method,
    short-plural=s,
    long-plural=s,
}

%%% CUSTOM COMMANDS
\newcommand{\ee}{\mathrm{e}}
\newcommand{\ii}{\mathrm{i}}
\newcommand{\pii}{\uppi}
\newcommand{\ed}{\downarrow}
\newcommand{\eu}{\uparrow}
\newcommand{\aB}{a_\mathrm{B}}
\newcommand{\me}{m_\mathrm{e}}
\newcommand{\EH}{E_\mathrm{H}}
\newcommand{\kBT}{k_\mathrm{B}T}
\newcommand{\alkcalc}{\texttt{AlkCalc} }
\newcommand{\figref}[2]{Fig.~\hyperref[#1]{\ref{#1}(#2)}}
\newcommand{\tus}{\rule{2mm}{.7pt}}

%%% TESTING
\usepackage{lipsum}

%%% TITLE PAGE
\title{Theory \& Manual\par\vspace*{5mm}for\par\vspace*{5mm}\alkcalc}%
\setcounter{footnote}{1}%
\author{Simon Euchner\footnote{Electronic Mail: %
        $\langle$\texttt{euchner.se@gmail.com}$\rangle$}%
       }
\date{20.08.2025}

%%% TABLES
\def\TabAngularFactors{
    \begin{table}[t]
        \centering
        \begin{tabular}{lllc}
            \toprule
            $(l,j)$ & $\mapsto$ & $(l',j')$ & $a_l$ \\
            \midrule
            $(l,l+s)$ & $\mapsto$ & $(l+1,l+s)$  & $\frac{1}{(2l+3)(2l+1)}$ \\
            $(l,l-s)$ & $\mapsto$ & $(l+1,l+s)$  & $\frac{l+1}{2l+1}$       \\
            $(l,l+s)$ & $\mapsto$ & $(l+1,l+3s)$ & $\frac{l+2}{2l+3}$       \\
            $(l,l-s)$ & $\mapsto$ & $(l-1,l-s)$  & $\frac{1}{(2l+1)(2l-1)}$ \\
            $(l,l+s)$ & $\mapsto$ & $(l-1,l-s)$  & $\frac{l}{2l+1}$         \\
            $(l,l-s)$ & $\mapsto$ & $(l-1,l-3s)$ & $\frac{l-1}{2l-1}$       \\
            \bottomrule
        \end{tabular}
        \caption{\textbf{Angular factors.} Table containing the angular factors
                 $a_l$ [see Eq.~\eqref{eq:OscillatorStrengthSimplified}] for
                 computing oscillator strengths. Throughout this table,
                 $s=1\slash{2}$.}
        \label{tab:AngularFactors}
    \end{table}
}



%%% START OF DOCUMENT
\begin{document}

    %%% TITLE PAGE
    \renewcommand{\pagemark}{}
    \maketitle%
    \newpage%

    %%% SET PLAIN PAGE STYLE
    \pagestyle{plain.scrheadings}%

    %%% TABLE OF CONTENTS

    % LINE FILLING
    \def\TOCLineFilling{\xleaders\hbox to.6em{\hfil{\tiny$\circ$\hfil}}\hfil}%

    % SET STYLES
    \DeclareTOCStyleEntry[linefill=\TOCLineFilling]{tocline}{section}%

    % PRINT
    \vspace*{2cm}
    \tableofcontents%
    \vfill
    {Typeset in \emph{Linux Libertine} with pdf-\LaTeXe \ in combination with
     \KOMAScript \ and \emph{latexmk}}
    \newpage

    %%% SET PAGE STYLE
    \pagestyle{scrheadings}%
    \setcounter{page}{1}%
    \renewcommand{\pagemark}{\thepage}%

    %%% SECTION
    \section{Introduction}
    \label{sec:introduction}

    Bound single-electron eigenfunctions of alkali-metal atoms and
    alkaline-earth metal ions are of significant relevance in quantum computing,
    quantum information processing, quantum simulation, Rydberg physics, quantum
    optics, \dots. However, expect for the special case of the Hydrogen atom, it
    is not possible to solve the time-independent Schr\"odinger equation
    analytically. Because of this, there is the need to obtain eigenenergies and
    eigenstates numerically. The library \alkcalc is designed for exactly this
    purpose: it allows the user to compute the relevant part of the bound-state
    energy spectrum with the associated eigenstates.

    The software \alkcalc is split into two parts: a first part used to compute
    the eigenenergies and eigenstates \emph{once} and store the resulting data.
    And a second part, referred to as the \emph{library functions}, which
    implement useful functionality for reading and processing the precomputed
    atomic data. For instance, there is as function to compute radial transition
    matrix elements.

    In the various \texttt{README} files included in the file tree of \alkcalc,
    technical details for computing the eigenenergies and eigenstates are
    provided. However, these files do not include a discussion on the theory and
    physics behind the code. Further, there is no reference manual for the
    library functions.

    This documents purpose is to add the missing pieces: in
    Sec.~\ref{sec:Theory} the theory and the physics associated with the code
    are discussed. Section \ref{sec:Manual} is a comprehensive reference manual
    for the library functions of \alkcalc.

    %%% SECTION
    \section{Theory}
    \label{sec:Theory}

    This section is organised as follows: in
    Subsec.~\ref{subsec:SingleAtomIonTheory} the theory and the physics of
    single-atom/ion systems are discussed. In
    Subsec.~\ref{subsec:NumericalMethod} the numerical method based on B-splines
    is presented. In Subsec.~\ref{subsec:RadialMatrixElements} it is shown how
    radial transitions matrix elements are computed by \alkcalc.  In
    Subsec.~\ref{subsec:OscillatorStrengths} it is shown how oscillator
    strengths for electronic transitions are computed. Finally, in
    Subsec.~\ref{subsec:OscillatorStrengths} it is shown how lifetimes of
    electronic states are computed by \alkcalc.

    %%% SUBSECTION
    \subsection{Single-atom/ion theory}
    \label{subsec:SingleAtomIonTheory}

    The starting point is an alkali-metal atom or alkaline-earth-metal ion. We
    are interested in systems where all electrons are in closed shells except
    for one valence electron. Examples of such systems are $\mathrm{^1H}$,
    $\mathrm{^{38}Rb}$, $\mathrm{^{88}Sr^+}$. This property allows a
    Kepler-problem description for the `effective nucleus' (the nucleus
    `screened' by the closed-shell electrons) and the valence electron. The
    dynamics of the valence electron and the effective nucleus can be described
    with the following Hamiltonian:
    \begin{equation}\label{eq:RawHamiltonian}
        H = \frac{\bm{p}^2}{2m}+V(r;\bm{L}^2,\bm{S}^2,\bm{J}^2)
    \,.
    \end{equation}
    Here, $m=Mm_\mathrm{e}\slash(M+m_\mathrm{e})$ is the reduced mass of the
    system, composed of the electron's mass $m_\mathrm{e}$ and the atom's or
    ion's mass $M$. The centre-of-mass motion is decoupled and $\bm{r},\bm{p}$
    are the relative canonical co-ordinates, i.e., $\bm{r}$ is the distance
    vector between the centre-of-mass position of the effective nucleus and the
    position of the valence electron. Further, $\bm{L}$, $\bm{S}$, and
    $\bm{J}=\bm{L}+\bm{S}$ are the orbital angular momentum operator, the
    electron's spin operator, and the total angular momentum operator,
    respectively. Finally, the radial potential $V$, depending on the distance
    $r=\norm*{\bm{r}}$, describes the interaction between the valence electron
    and the effective nucleus. It accounts for several effects: far from the
    nucleus the valence electron experiences a Coulomb potential with effective
    charge $Z_\mathrm{eff}=1$ for alkali-metal atom and $Z_\mathrm{eff}=2$ for
    singly ionised alkaline-earth metal atoms.  This is because the nucleus of
    charge $Z$ is `screened' by the $Z-1$ (or $Z-2$, for alkaline-earth metal
    ions) ground-state electrons in closed-shells. However, valence electrons
    with orbits sufficiently close to the nucleus experience the non-trivial
    charge distribution posed by the nucleus in conjugation with the tightly
    bound electron cloud. This leads to short-ranged corrections to the
    effective Coulomb potential with charge $Z_\mathrm{eff}$ at long distances.
    Following Refs.~\cite{aymar1996,marinescu1994,condon1951,wilkinson2025}
    these corrections are described by the following potential:
    \begin{equation}\label{eq:ParametricModelPotential}
        V(r;\bm{L}^2,\bm{S}^2,\bm{J}^2)
        =
        V_\mathrm{C}(r)
        +
        V_\mathrm{P}(r)
        +
        V_\mathrm{R}(r;\bm{L}^2,\bm{S}^2,\bm{J}^2)
    \,.
    \end{equation}
    The first contribution is
    \begin{equation}
        V_\mathrm{C}(r)
        =
        -\frac{e^2}{4\pii\varepsilon_0}\frac{Z_\mathrm{n}(r)}{r}
    \,, \
        Z_\mathrm{n}(r)
        =
        Z_\mathrm{c}+(Z-Z_\mathrm{c})\ee^{-k_1r}
        +
        k_2r\ee^{-k_3r}
        -
        k_4r^2\ee^{-k_3r}
    \,.
    \end{equation}
    This is a modified Coulomb potential which has the tail of a Coulomb
    potential with charge $Z_\mathrm{eff}$ and short-ranged corrections for
    small radii accounting for the charge-distribution of the effective nucleus.
    Here, the parameters $k_j$ ($j=1,2,3,4$) are empirically determined from
    spectroscopic data \cite{aymar1996,marinescu1994}. The second contribution
    to the potential \eqref{eq:ParametricModelPotential} is
    \begin{equation}
        V_\mathrm{P}(r)
        =
        -\frac{e^2}{(4\pii\varepsilon_0)^2}
        \frac{\alpha_\mathrm{d}}{2r^4}
        \bigg[1-\exp(-\frac{r^6}{r^6_\mathrm{c}})\bigg]
    \,.
    \end{equation}
    It accounts for the polarisability of the screened nucleus, which gets
    polarised by the presence of the valence electron. Here, $r_\mathrm{c}$ is
    the effective nucleus' size and $\alpha_\mathrm{d}$ its static electric
    polarisability. Finally,
    \begin{equation}\label{eq:LSCoupling}
        V_\mathrm{R}(r;\bm{L}^2,\bm{S}^2,\bm{J}^2)
        =
        \frac{1}{2m^2c^2}\frac{V'_\mathrm{NR}(r)}{rN(r)}\bm{L}\cdot\bm{S}
    \,, \
        N(r) = \bigg[ 1 - \frac{V_\mathrm{NR}(r)}{2mc^2} \bigg]^2
    \,, \
        V_\mathrm{NR}(r) = V_\mathrm{C}(r) + V_\mathrm{P}(r)
    \,.
    \end{equation}
    It describes (relativistic) coupling between the orbital angular momentum
    and the valence electrons' spin. This expression results from relativistic
    perturbation theory \cite{condon1951}. Here, $V_\mathrm{NR}$ is the
    non-relativistic model potential and $c$ the speed of light. In
    Ref.~\cite{wilkinson2025} the potential \eqref{eq:LSCoupling} is employed
    for the treatment of alkaline-earth metal ions. In Ref.~\cite{sibalic2017}
    the simplified version
    \begin{equation}\label{eq:SimplifiedLSCoupling}
        V_\mathrm{R}(r)
        =
        \frac{\alpha}{2r^3} \bm{L}\cdot\bm{S}
    \end{equation}
    is employed for alkali-metal atoms. Here, $\alpha$ is the fine-structure
    constant. Throughout \alkcalc relativistically correct version
    in Eq.~\eqref{eq:LSCoupling}, not the simplified version in
    Eq.~\ref{eq:SimplifiedLSCoupling}, is employed. Finally, for expressing the
    spin-operator $\bm{S}$ in Eq.~\eqref{eq:LSCoupling} and
    Eq.~\eqref{eq:SimplifiedLSCoupling} the representation in Pauli's (spin)
    matrices is chosen:
    \begin{equation}
        \bm{S} = \frac{\hbar}{2}\mqty[\sigma_x \\ \sigma_y \\ \sigma_z]
    \,.
    \end{equation}
    The Pauli spin matrices are
    \begin{equation}
        \sigma_x = \mqty[0&1 \\ 1&0]
    \,, \
        \sigma_y = \mqty[0&-\ii \\ \ii&0]
    \,, \
        \sigma_z = \mqty[1&0 \\ 0&-1]
    \,.
    \end{equation}
    This simple choice for representing the spin degreed of freedom is possible
    because throughout this work the interest is in the case where the total
    spin is $s=1\slash{2}$.

    At this point the Hamiltonian which determines the eigenenergies and the
    eigenstates is defined. In the following, the goal is to cast eigenvalue
    problem posed by Hamiltonian \eqref{eq:RawHamiltonian} into a more practical
    form for the numerical procedure. The first step is to rewrite the product
    $\bm{L}\cdot\bm{S}$ in terms of the total angular momentum operator
    $\bm{J}=\bm{L}+\bm{S}$:
    \begin{equation}
        \bm{L}\cdot\bm{S}
        =
        \frac{1}{2}\big[\bm{J}^2-\bm{L}^2-\bm{S}^2\big]
    \,.
    \end{equation}
    Next, the square of the momentum operator in spherical co-ordinates is
    rewritten:
    \begin{equation}
        \bm{p}^2
        =
        (-\ii\hbar\nabla)^2
        =
        -\hbar^2\nabla^2
        =
        -\frac{\hbar^2}{r^2}\pdv{r}(r^2\pdv{r})+\frac{\bm{L}^2}{r^2}
    \,.
    \end{equation}
    The second contribution is the centrifugal barrier. In a classical Kepler
    problem this is what prevents the electron crashing into the nucleus.
    Expressed in the spherical coordinates, Hamiltonian
    \eqref{eq:RawHamiltonian} takes the form
    \begin{equation}
        H
        =
        -\frac{\hbar^2}{2mr^2}\pdv{r}(r^2\pdv{r})
        +
        V_\mathrm{eff}(r;\bm{L}^2,\bm{S}^2,\bm{J}^2)
    \,, \
        V_\mathrm{eff}(r;\bm{L}^2,\bm{S}^2,\bm{J}^2)
        =
        V(r;\bm{L}^2,\bm{S}^2,\bm{J}^2)
        +
        \frac{\bm{L}^2}{2mr^2}
    \,.
    \end{equation}
    The full Hilbert space this Hamiltonian acts on can be split into an angular
    momentum factor and a radial one. A basis of the angular momentum Hilbert
    space is given by the simultaneous eigenstates $\Phi_{l,s,j,m_j}$ of the
    operators (Cartan subalgebra) $\bm{L}^2$, $\bm{S}^2$, $\bm{J}^2$, and $J_z$.
    Note that a simultaneous eigenbasis of $\bm{L}^2$, $\bm{S}^2$, $\bm{J}^2$,
    and $J_z$ exists, because these operators commute pairwise. The so-called
    \emph{spinors} $\bm{\Phi}_{l,s,j,m_j}$ constitute the so-called
    \emph{coupled basis}, also known as the \emph{fine-structure} basis in the
    literature. These spinors are given by linear combinations of the uncoupled
    basis states
    \begin{equation}\label{eq:UncoupledSpinors}
        Y_{l,m_l}\otimes\bm{\chi}_{s,m_s}
    \end{equation}
    via SU(2) Clebsch-Gordan coefficients \cite{friedrich2017}. Here,
    $Y_{l,m_l}\colon[0,\pii]\times[0,2\pii]\to\mathbb{C}$ are spherical
    harmonics defined as
    \begin{equation}\label{eq:sphericalHarmonics}
        Y_{l,m_l}(\vartheta,\varphi)
        =
        \sqrt{\frac{2l+1}{4\pii}\frac{(l-m_l)!}{(l+m_l)!}}
        P^{m_l}_l(\cos(\vartheta))\,\ee^{\ii m_l \varphi}
    \,.
    \end{equation}
    In this expression $P^{m_l}_l\colon[-1,1]\to\mathbb{R}$ are the associated
    Legendre polynomials defined for $l=0,1,2,\dots$ and $m_l=0,1,\dots,l-1,l$.
    These polynomials can be computed via the recursion
    \begin{equation}
        (l-m_l) P^{m_l}_l(x)
        =
        x (2l-1) P^{m_l}_{l-1}(x)
        -
        (l+m_l-1) P^{m_l}_{l-2}(x)
    \,,
    \end{equation}
    with starting point
    \begin{equation}
        P^{m_l}_{m_l}(x)
        =
        (-1)^{m_l}\frac{(2m_l)!}{2^{m_l}m_l!} (1-x^2)^{m_l\slash{2}}
    \,,
    \end{equation}
    in combination with $P^{m_l}_l(x)=0$ for $l<m_l$. Negative values of $m_l$
    are treated with the relation
    \begin{equation}
        P^{-m_l}_l = (-1)^{m_l} \frac{(l-m_l)!}{(l+m_l)!}P^{m_l}_l
    \,.
    \end{equation}
    In this definition of the associated Legendre polynomial the so-called
    \emph{Condon-Shortley phase} is included. This follows the definition in
    Ref.~\cite{friedrich2017}.\footnote{This definition entails that the factor
    $(-1)^{m_l}$ does not explicitly appear in the definition of the spherical
    harmonics. This remark is placed here to avoid confusion of whether the
    spherical harmonics are with or without the Condon-Shortley phase. In the
    definition of $Y_{l,m_l}$ used in this work, the Condon-Shortley phase
    included.} Finally, $\bm{\chi}_{s,m_s}$ in Eq.~\eqref{eq:UncoupledSpinors}
    are two-spinors defined as
    \begin{equation}\label{eq:twoSpinsors}
        \bm{\chi}_{\frac{1}{2},\frac{1}{2}}=\mqty[1 \\ 0]
    \,, \
        \bm{\chi}_{\frac{1}{2},-\frac{1}{2}}=\mqty[0 \\ 1]
    \,.
    \end{equation}
    Explicitly (for the case of interest $s=1\slash{2}$) the coupled basis
    states $\bm{\Phi}_{l,s,j,m_j}$ are:
    \begin{equation}\label{eq:coupledBasis}
        \bm{\Phi}_{l,s,j,m_j}
        =
        \langle Y_{l,m_j-s}\otimes\bm{\chi}_{s,s},\bm{\Phi}_{l,s,j,m_j} \rangle
        Y_{l,m_j-s}\otimes\bm{\chi}_{s,s}
        +
        \langle Y_{l,m_j+s}\otimes\bm{\chi}_{s,-s},\bm{\Phi}_{l,s,j,m_j} \rangle
        Y_{l,m_j+s}\otimes\bm{\chi}_{s,-s}
    \,,
    \end{equation}
    where the scalar products are (possibly vanishing) Clebsch-Gordan
    coefficients. These spinor transform according to irreducible
    representations of SU(2). This is seen in the following way: consider the
    unitary and faithful representation
    $\rho(\bm{\alpha})=\exp(\ii\bm{\alpha}\cdot\bm{J})$, where
    $\bm{\alpha}\in\mathbb{R}^3$ is a vector of length
    $\norm*{\bm{\alpha}}\in[0,2\pii)$. Its action on the spinors is
    \begin{equation}
        \rho(\bm{\alpha})
        \Phi_{l,s,j,m_j}(\vartheta,\varphi)
        =
        \sum^{d_j}_{q=1} D^{(\Gamma_j)}(\bm{\alpha})_{m_j,q}
        \Phi_{l,s,j,q}(\vartheta,\varphi)
    \,.
    \end{equation}
    In this expression, $d_j$ is the dimension of the irreducible representation
    $\Gamma_j$ of SU(2) and $D^{(\Gamma_j)}$ is a unitary matrix representation
    of the irreducible representation $\Gamma_j$.

    To rewrite the Hamiltonian the crucial step is to use the fact that it has
    SU(2) symmetry, i.e., $\comm*{\rho(\bm{\alpha})}{H}=0$ for all
    $\bm{\alpha}$.\footnote{The symmetry group is even bigger because parity is
    also conserved. So the full symmetry group is
    $\mathrm{SU}(2)\times\mathbb{Z}_2$.} Due to this symmetry the Hamiltonian is
    automatically block diagonal with block dimensions given by integer
    multiplicities $\mu_j$ of the dimensions $d_j=2j+1$ of the irreducible
    representations $\Gamma_j$ labelled by
    $j=0,1\slash{2},1,3\slash{2},2,\dots$. The multiplicities $\mu_j$ are
    $\mu_{j}=2$ for $j\neq0$ and $\mu_0=0$. This is because for each $l>0$ there
    are the two possible orbital angular momentum quantum numbers
    $l=j-1\slash{2}$ and $l=j+1\slash{2}$ and for $s=1\slash{2}$ a total angular
    momentum of $j=0$ is not possible. This entails that the Hamiltonian in the
    basis of the spinors $\bm{\Phi}_{l,s,j,m_j}$ can be decomposed into the
    following direct sum:
    \begin{equation}
        \bigoplus^\infty_{l=0}
        \bigoplus_{j \in \{\abs*{l-s},l+s\}}
        \bigoplus^j_{m_j=-j}
        H_{l,s,j,m_j}
    \,, \ \ \
        H_{l,s,j,m_j}
        =
        \langle \bm{\Phi}_{l,s,j,m_j}, H \bm{\Phi}_{l,s,j,m_j} \rangle
    \,.
    \end{equation}
    Note that the off-diagonal matrix elements within the blocks associated with
    each irreducible representation $\Gamma_j$ vanish, i.e.,
    $\langle\bm{\Phi}_{l,s,j,m'_j},H\bm{\Phi}_{l,s,j,m_j}\rangle=0$ for
    $m_j'\neq{m}_j$. This is a result of the
    $\mathrm{U}(1)\subset\mathrm{SU}(2)$ symmetry of the full Hamiltonian in
    Eq.~\eqref{eq:RawHamiltonian}, which is associated to the generator $L_z$.
    This symmetry further entails that the diagonal matrix elements
    $\langle\bm{\Phi}_{l,s,j,m_j},H\bm{\Phi}_{l,s,j,m_j}\rangle$ do not depend
    on $m_j$. Therefore, the full Hamiltonian \eqref{eq:RawHamiltonian} in the
    basis of the spinors $\bm{\Phi}_{l,s,j}$ can be expressed as the more simple
    direct sum
    \begin{equation}
        \bigoplus^\infty_{l=0}
        \bigoplus_{j \in \{\abs*{l-s},l+s\}}
        1_{2j+1} \otimes H_{l,s,j}
    \,.
    \end{equation}
    Here, $1_{2j+1}$ is the identity matrix of dimension $2j+1$ and
    \begin{equation}
        H_{l,s,j}
        =
        -\frac{\hbar^2}{2mr^2}\pdv{r}(r^2\pdv{r})
        +
        V_\mathrm{eff}(r;\hbar^2l[l+1],\hbar^2s[s+1],\hbar^2j[j+1])
    \,.
    \end{equation}
    For completeness $\rho$ can be expressed in the basis of
    $\bm{\Phi}_{l,s,j}$ as well. It is given by the following direct sum of
    the irreducible representations $\Gamma_j$ os SU(2):
    \begin{equation}
        \rho
        =
        \bigoplus_{j>0}
        2\Gamma_j
    \,.
    \end{equation}
    The block structure of the Hamiltonian reduces the eigenvalue problem
    defined by the time-independent Schr\"odinger equation
    \begin{equation}
        H \bm{\psi} = E \bm{\Psi}
    \end{equation}
    to the blocks $H_{l,s,j}$ on the Hilbert space associated with the radial
    degrees of freedom. In particular, solving for the full eigenenergies and
    eigenstates is reduced to solving the radial eigenvalue problems
    \begin{equation}\label{eq:radialEigenproblem}
        H_{l,s,j} R_{n,l,s,j} = E_{n,l,s,j} R_{n,l,s,j}
    \,,
    \end{equation}
    where $R_{n,l,s,j}$ is the $n$-th radial eigenstate of $H_{l,s,j}$. The
    total eigenstate $\bm{\Psi}$ is then given by the product of the radial
    eigenstates and the spinors:
    \begin{equation}\label{eq:GeneralRadialEquations}
        \bm{\Psi}_{n,l,s,j,m_j}(r,\vartheta,\varphi)
        =
        R_{n,l,s,j}(r)
        \bm{\Phi}_{l,s,j,m_j}(\vartheta,\varphi)
    \,.
    \end{equation}
    This shows that solving for the full eigenstates $\bm{\Psi}_{n,l,s,j,m_j}$
    is equivalent to solving for the radial eigenstates $R_{n,l,s,j}$.  In
    general, the eigenenergies $E=E_{n,l,s,j}$ and the radial eigenstates
    $R_{n,l,s,j}$ must be calculated numerically. This task is exactly what
    \alkcalc is designed to do.

    {\color{blue}
    %%% SUBSECTION
    \subsection{Numerical method}
    \label{subsec:NumericalMethod}

    In this subsection the numerical method which is employed by \alkcalc to
    solve the radial eigenvalue problems in
    Eq.~\eqref{eq:GeneralRadialEquations} for the eigenenergies and the radial
    eigenstates $R_{n,l,s,j}$ is discussed.  The first step is to express the
    eigenvalue problems in terms of dimensionless quantities. To this end the
    Hamiltonians $H_{l,s,j}$ is expressed in units of Hartree,
    $E_\mathrm{H}=e^2\slash{(4\pii\varepsilon_0a_\mathrm{B})}$, where
    $a_\mathrm{B}=\hbar\slash{(m_\mathrm{e}c\alpha)}\approx52.9\,\mathrm{pm}$ is
    Bohr's radius. Further, the radial eigenfunctions are parametrised as
    \begin{equation}\label{eq:DimensionlessEigenfunctions}
        R_{n,l,s,j}(r)
        =
        \frac{1}{\sqrt{\aB^3}}\frac{f_{n,l,s,j}(r\slash\aB)}{r\slash\aB}
        \,,
    \end{equation}
    where $f_{n,l,s,j}$ is dimensionless. The eigenvalue problems in
    Eq.~\eqref{eq:GeneralRadialEquations} expressed in the dimensionless
    eigenfunctions \eqref{eq:DimensionlessEigenfunctions}
    takes the form
    \begin{equation}\label{eq:Eigenproblem}
        \bigg[
        -
        \frac{1}{2}
        \dv[2]{t}
        +\frac{l(l+1)}{2t^2}
        +C\frac{V(\aB{t};\hbar^2l[l+1],\hbar^2s[s+1],\hbar^2j[j+1])}{\EH}
        \bigg]
        f_{n,l,s,j}(t)
        =
        \bar{E}_{n,l,s,j}
        f_{n,l,s,j}(t)
    \,,
    \end{equation}
    where
    \begin{equation}
        C = \frac{1}{1+\me\slash{M}}
    \end{equation}
    and the eigenenergies may be extracted from the relation
    \begin{equation}
        \frac{E_{l,s,j,n}}{E_\mathrm{H}}
        =
        \frac{\bar{E}_{n,l,s,j}}{C}
    \,.
    \end{equation}
    Rewriting the normalisation condition in terms of the eigenfunctions
    $f_{n,l,s,j}$ leads to
    \begin{equation}\label{eq:normalisation}
        1
        =
        \int^\infty_0\dd{r}r^2R^2_{n,l,s,j}(r)
        =
        \int^\infty_0 (\aB\dd{t})(\aB t)^2 \frac{1}{\aB^3}
        \frac{f^2_{n,l,s,j}(t)}{t^2}
        =
        \int^\infty_0\dd{t}f^2_{n,l,s,j}(t)
    \,.
    \end{equation}
    Here is is assumed that the radial eigenfunctions are real-valued, which is
    justified as the Hamiltonian is not only an hermitian but also a symmetric
    operator. In conclusion, computing the normalised eigenfunctions
    $f_{n,l,s,j}$ yields the (automatically normalised) radial eigenstates by
    employing Eq.~\eqref{eq:DimensionlessEigenfunctions} for the conversion.

    To solve the eigenvalue problem in Eq.~\eqref{eq:Eigenproblem} numerically,
    so-called \emph{B-splines} are employed.\footnote{This method may be
    conceptually and programmatically more involved than a simple
    finite-difference or finite-element method, however, it is preferred here
    for two reasons: it allows to work with much smaller matrix dimensions and
    it allows to store the resulting eigenenergies and radial eigenstates much
    more efficiently.} Ultimately, this allows to express the eigenvalue problem
    for the radial eigenfunctions $f_{n,l,s,j}$ as a simple matrix eigenvalue
    problem. The remainder of this subsection will be concerned with
    two different steps: first, the notion of B-splines is introduced. Second,
    the final form of the matrix eigenproblem which yields the eigenenergies and
    the radial eigenstates is formulated.

    The first step in setting-up the final matrix eigenvalue problem is to agree
    on a finite closed interval $I=[0, t_\mathrm{tmax}]$, where
    $t_\mathrm{max}>0$ sets the maximal value of the domain we consider. The
    value $t_\mathrm{max}$ must be chosen large enough to ensure that the
    boundary condition $f_{n,l,s,j}(t\geq{t}_\mathrm{max})=0$ is satisfied. This
    equal sign is actually never true, because the radial eigenstates
    asymptotically approach zero. However, if zero is reached within machine
    precision, the boundary condition is numerically satisfied. Of course, one
    must ensure that the finite set of solutions that shall be computed has
    support within this interval. The mesh we define as a family
    $(t_k)_{k=1,\dots,N-1}$ in $I$ with the properties $t_0=0$,
    $t_{N-1}=t_\mathrm{max}$, and $t_0<t_1<\dots<t_{N-2}<t_{N-1}$, where we say
    $N\geq{3}$. Of course, typical values are $N\,{\sim}\,10^5\text{-}10^7$.
    Next, we define the derived family of step sizes $(h_k)_{k=1,\dots,N-1}$,
    $h_k=t_k-t_{k-1}$. The elements $(\phi_k)_{k=0,\dots,N-1}$,
    $\phi_k\colon{I}\to[0,1]$, we define as
    \begin{equation}
        \phi_k(t)
        =
        \begin{cases}
            (t-t_{k-1})\slash(t_k-t_{k-1})\,, \ & t\in(t_{k-1},t_k)
        \\
            (t_{k+1}-t)\slash(t_{k+1}-t_k)\,, \ & t\in(t_k,t_{k+1})
        \\  0                             \,, \ & \text{else}
        \end{cases}
    \,.
    \end{equation}
    for $k=1,\dots,N-2$, and for $k=0$ and $k=N-1$ the Dirichlet boundary
    conditions $f_{n,l,s,j}(0)=f_{n,l,s,j}(t_\mathrm{max})=0$ lead to the
    canonical choice $\phi_0=\phi_{N-1}\equiv0$. To obtain a matrix problem
    which approximates Eq.~\eqref{eq:eigenproblem} we compute projections of it
    onto the elements with the real scalar product (`dummy' functions
    $\alpha,\beta\colon{I}\to\mathbb{R}$)
    \begin{equation}
        B(\alpha,\beta)
        =
        \int_I \dd{t} \alpha(t) \beta(t)
    \,.
    \end{equation}
    To simplify the notation we fix the quantum numbers $n$, $l$, $s$, and $j$
    and say $f\equiv{f_{n,l,s,j}}$. Moreover, we collect all potential
    contributions into the new potential $v$, defined as
    \begin{equation}
        v(t)
        =
        \frac{l(l+1)}{2t^2}
        +C\frac{V(\aB{t};\hbar^2l[l+1],\hbar^2s[s+1],\hbar^2j[j+1])}{\EH}
    \,.
    \end{equation}%
    %
    Finally, we say $\lambda=\bar{E}_{n,l,s,j}$, such that
    Eq.~\eqref{eq:eigenproblem} simply becomes
    \begin{equation}
        -\frac{1}{2}f''(t) + v(t) f(t) = \lambda f(t)
    \,.
    \end{equation}
    Before we project this equation onto the elements, we approximately expand
    $f$ in them:
    \begin{equation}\label{eq:expansionOfState}
        f = \sum^{N-2}_{k=1} f_k \phi_k
    \,.
    \end{equation}
    This expansion automatically fulfils the Dirichlet boundary conditions
    $f(0)=f(t_\mathrm{max})=0$. Projection, i.e., applying the map
    $B(\phi_k,\cdot)$ for all $k=1,\dots,N-2$, yields a matrix eigenproblem
    whose solutions converge to solutions of Eq.~\eqref{eq:eigenproblem} as $N$
    tends to infinity:
    \begin{equation}\label{eq:matrixEigenproblem}
        (K+V) \bm{f} = \lambda M \bm{f}
    \,.
    \end{equation}
    In this expression $M_{kk'}=B(\phi_k,\phi_{k'})$ is the so-called mass
    matrix, $K_{kk'}=B(\phi'_k,\phi'_{k'})\slash{2}$ the so-called stiffness
    matrix\footnote{Note that because the elements obey the Dirichlet boundary
    conditions $\phi_k(0)=\phi_k(t_\mathrm{max})=0$, $k=1,\dots,N-2$,
    integrating by parts allows to write
    $B(\phi''_k,\phi_{k'})=-B(\phi'_k,\phi'_{k'})$.}, and
    $V_{kk'}=B(\phi_k,v\phi_{k'})$ the potential energy contribution. Each
    matrix is $(N-2)$-dimensional, since $k,k'=1,\dots,N-2$, and $M$ and $K$ can
    be computed analytically; in particular
    \begin{equation}
        M
        =
        \frac{1}{6}
        \mqty[2(h_1+h_2) &  h_2 & & & \\
              h_2 & 2(h_2+h_3) & h_3 & & \\
              & h_3 & 2(h_3+h_4) & h_4 & \\
              & & \ddots & \ddots & h_{N-2} \\
              & & & h_{N-2} & 2(h_{N-2}+h_{N-1})]
    \end{equation}
    and
    \begin{equation}
        K = \frac{1}{2}
            \mqty[\frac{1}{h_1}+\frac{1}{h_2} & -\frac{1}{h_2} & \\
                  -\frac{1}{h_2} & \frac{1}{h_2}+\frac{1}{h_3} & -\frac{1}{h_3}
                  & \\
                  & -\frac{1}{h_3} & \frac{1}{h_3}+\frac{1}{h_4} &
                  -\frac{1}{h_4} & \\
                  & & \ddots & \ddots & -\frac{1}{h_{N-2}} \\
                  & & & -\frac{1}{h_{N-2}} & \frac{1}{h_{N-2}}+\frac{1}{h_{N-1}}
                  ]
    \,.
    \end{equation}
    Note that both of these matrices are symmetric. The potential matrix,
    generally, must be computed numerically. However, here we choose the simple
    approximation
    \begin{equation}
        B(\phi_k,v\phi_{k'})
        \approx
        \frac{v(t_k)+v(t_{k'})}{2}
        M_{k{k'}}
    \,,
    \end{equation}
    which is justified if the step size is small enough. Note that this
    particular choice of approximations keeps the matrix symmetric, such that
    $K+V$ is symmetric. Furthermore, $M$ represents a scalar product and is thus
    positive definite. This means that the matrix eigenproblem in
    Eq.~\eqref{eq:matrixEigenproblem} can be efficiently solved with the Lanczos
    algorithm.

    Next, we express the normalisation condition in
    Eq.~\eqref{eq:normalisation} in terms of our elements and the vector
    $\bm{f}$. To this end we calculate
    \begin{equation}
        \int^\infty_0\dd{t}f^2(t)
        \approx
        \int_I  \dd{t} f^2(t)
        =
        \sum^{N-2}_{k=1}\sum^{N-2}_{k'=1}f_kf_{k'} B(\phi_k,\phi_{k'})
        =
        \bm{f}^\mathrm{T}M\bm{f}
        =
        \norm*{\bm{f}}^2_M
        \overset{!}{=}
        1
    \,,
    \end{equation}
    i.e., the length of $\bm{f}$ must be unity in the by the scalar product $M$
    induced norm $\norm*{\,\cdot\,}_M$.

    %%% SUBSECTION
    \subsection{Radial matrix elements}
    \label{subsec:RadialMatrixElements}

    To compute radial matrix elements, we first calculate, for $p\in\mathbb{R}$
    such that the integral exists,
    \begin{align}\label{eq:radialMatrixElement}
        \begin{split}
            \frac{r^{n,l,s,j;n',l',s',j'}_p}{\aB^p}
            &\equiv
            \frac{1}{\aB^p} \langle R_{n,l,s,j},
            \mathrm{id}^p R_{n',l',s',j'} \rangle
            =
            \int^\infty_0 \dd{r} r^2
            R_{n,l,s,j}(r)
            \frac{r^p}{\aB^p}
            R_{n',l',s',j'}(r)
        \\
            &=
            \int^\infty_0 (\dd{t} \aB) (\aB t)^2
            f_{n,l,s,j}(t)
            \frac{(\aB t)^p}{\aB^p} \frac{1}{\aB^3}
            \frac{f_{n',l',s',j'}(t)}{t^2}
            =
            B(f_{n,l,s,j}, \mathrm{id}^p f_{n',l',s',j'})
        \,.
        \end{split}
    \end{align}
    Here, `$\mathrm{id}$' denotes the identity on the non-negative real numbers.
    Next, we approximate this result with solutions to the matrix eigenproblem
    in Eq.~\eqref{eq:matrixEigenproblem}, i.e., we represent $f_{n,l,s,j}$ and
    $f_{n',l',s',j'}$ approximately with the vectors $\bm{f}_{n,l,s,j}$ and
    $\bm{f}_{n',l',s',j'}$, respectively [see Eq.~\eqref{eq:expansionOfState}].
    This allows to write
    \begin{equation}
        B(f_{n,l,s,j}, \mathrm{id}^p f_{n',l',s',j'})
        \approx
        \bm{f}_{n,l,s,j}^\mathrm{T} R_p \bm{f}_{n',l',s',j'}
    \,,
    \end{equation}
    where the components of the real symmetric matrix $R_p$ are defined as
    \begin{equation}
        (R_p)_{k,k'}
        =
        B(\phi_k,\mathrm{id}^p\phi_{k'})
        =
        \begin{cases}
            B(\phi_k,\mathrm{id}^p\phi_{k+1}),
            & k'=k+1 \ \text{or} \ k=k'+1 \ (0<k,k'<N-2)
        \\
            B(\phi_k,\mathrm{id}^p\phi_{k}),
            & k'=k \ (0<k,k'\leq{N-2})
        \\
        0, & \text{else}
        \end{cases}
    \,.
    \end{equation}
    To numerically compute the remaining integrals we use a two-point
    Gau\ss{i}an quadrature. This method approximates the integral as
    \cite{burden2011}:
    \begin{equation}
        \int^b_a \dd{t} \alpha(t)
        \approx
        \frac{b-a}{2}
        \bigg[
            \alpha\bigg(\frac{a+b}{2}-\frac{b-a}{2\sqrt{3}}\bigg)
            +
            \alpha\bigg(\frac{a+b}{2}+\frac{b-a}{2\sqrt{3}}\bigg)
        \bigg]
    \,,
    \end{equation}
    for any sufficiently smooth function $\alpha$ and $a<b$. Note that this
    result is exact if $\alpha$ is a polynomial of degree three, which is the
    case for radial dipole matrix elements. This method is straight-forward to
    implement in a computer program, such that the problem of calculating matrix
    elements is effectively reduced to computing matrix products. However, note
    that for too large $p$ there can be numerical overflow or underflow when
    computing $t^p$. But, for 64-bit floating point arithmetic, $p=\pm10$ is
    easily possible (if the integral exists), such that, typically, all
    practical cases are included. Explicitly, we could verify that for Hydrogen
    the range $p=-2,\dots,10$ yields reasonably good results. Finally, we remark
    that one can improve on the precision by employing higher-order Gau\ss{i}an
    quadrature algorithms as described in \cite{burden2011}.
    }

    %%% SUBSECTION
    \subsection{Oscillator strengths}
    \label{subsec:OscillatorStrengths}

    The so-called \emph{oscillator strength} of a transition from an initial
    state $\bm{\Psi}_\mathrm{i}\equiv\bm{\Psi}_{n',l',s,j',m'_j}$ to a final
    state $\bm{\Psi}_\mathrm{f}\equiv\bm{\Psi}_{n,l,s,j,m_j}$ is
    \cite{friedrich2017}
    \begin{equation}
        f_{\mathrm{i}\to\mathrm{f}}
        =
        \sum^{j'}_{m'_j=-j'}
        \frac{2}{3} \bar{E}_\mathrm{f,i}
        \abs*{\langle\bm{\Psi}_\mathrm{f},
                     \bm{r}\slash\aB\bm{\Psi}_\mathrm{i}\rangle}^2
    \,,
    \end{equation}
    with the energy difference
    $\EH\bar{E}_\mathrm{f,i}=E_\mathrm{f}-E_\mathrm{i}$. Note that the sum runs
    over all final $m'_j$ because one is interested in transitions between the
    fine-structure states without selecting Zeeman sub-levels. The oscillator
    strength is a measure for how strongly the induced Hertzian (electric)
    dipole oscillates: the radiation power is large for spatially large dipoles
    and fast oscillation frequencies, captured by the matrix element and energy
    difference, respectively. In principle, all the machinery to numerically
    compute this quantity is already there, however, one can reduce the
    expression to a much simpler form by using angular momentum algebra, in
    particular, relations from \cite{varshalovic1988}. To simplify the
    expression, first note (Condon-Shortley phase convention)
    \begin{equation}
        r_x
        =
        \sqrt{\frac{2\pii}{3}} r (Y_{1,-1}-Y_{1,1})
    \,, \
        r_y
        =
        \ii \sqrt{\frac{2\pii}{3}} r (Y_{1,-1}+Y_{1,1})
    \,, \
        r_z
        =
        \sqrt{\frac{4\pii}{3}} r Y_{1,0}
    \,.
    \end{equation}
    With this, the oscillator strength becomes
    \begin{equation}
        f_{\mathrm{i}\to\mathrm{f}}
        =
        \frac{4\pii}{3} \frac{2}{3} \bar{E}_\mathrm{f,i}
        \langle{r}R_\mathrm{f},r^2\slash\aB{r}R_\mathrm{i}\rangle^2
        \sum^{j'}_{m'_j=-j'}
        \sum^1_{q=-1}
        \abs*{\langle\bm{\Phi}_\mathrm{f},Y_{1,q}\bm{\Phi}_\mathrm{i}\rangle}^2
    \end{equation}
    Employing relations from \cite[p.~236,244,260,445,496]{varshalovic1988} one
    can show that the remaining matrix element is
    \begin{equation}
        \langle \bm{\Phi}_\mathrm{f}, Y_{1,q}\bm{\Phi}_\mathrm{i} \rangle
        =
        (-1)^{j'+j+1+l'+s-m'_j}
        \sqrt{\frac{3}{4\pii}}
        \mqty(j' & 1 & j \\ -m'_j & q & m_j)
        \sqrt{(2j+1)(2j'+1)(2l+1)}
        \bigg\{\mqty{l & s & j \\ j' & 1 & l'}\bigg\}
        C^{l,0,1,0}_{l',0}
    \,,
    \end{equation}
    expressed in terms of a Wigner $3jm$-symbol, a Wigner $6j$-symbol, and a
    Clebsch-Gordan coefficient. With symmetry relations and a sum rule from
    \cite[p.~236,259]{varshalovic1988} one can compute the sum over $q$ and
    $m'_j$ in the expression for the oscillator strength, which yields
    \begin{equation}\label{eq:OscillatorStrengthSimplified}
        \sum_{m'_j} f_{\mathrm{i}\to\mathrm{f}}
        =
        \frac{2}{3} \bar{E}_\mathrm{f,i}
        \abs*{\langle r R_{n',l',s,j'}, r^2\slash\aB R_{n,l,s,j} \rangle}^2
        a_l
    \,, \
        a_l
        =
        l_\mathrm{max} (2j'+1)
        \bigg\{\mqty{l & s & j \\ j' & 1 & l'}\bigg\}^2
    \,,
    \end{equation}
    where $l_\mathrm{max}=\mathrm{max}(l,l')$. There are only six transitions
    for which $a_l$ is non-zero. Using relations from
    \cite[p.~298,300]{varshalovic1988} one arrives at the results shown in
    Tab.~\ref{tab:AngularFactors}. With this it is now straight forward to
    compute the oscillator strengths numerically with a computer program. As a
    final remark, note that the oscillator strength does not depend on $m_j$ any
    more, i.e., after summing over final states, the oscillators for different
    $m_j$ all have the same strength.

    %%% TABLE
    \TabAngularFactors

    %%% SUBSECTION
    \subsection{State lifetimes}
    \label{subsec:StateLifetimes}

    The lifetime $\tau$ of a fine-structure state is given by the formula
    \begin{equation}
        \tau = \Gamma^{-1}
    \,,
    \end{equation}
    where $\Gamma$ is the total rate at which the state is depopulated. At
    sufficiently low temperatures this rate is given by the sum over the
    so-called \emph{Einstein-A} coefficients $A_\mathrm{i\to{f}}$
    \cite{einstein1917,friedrich2017}. These coefficients are the rates at which
    the state of interest, here referred to as the `initial state', radiatively
    decays to some final state. In the dipole approximation the sum over
    Einstein-A coefficients can be expressed as \cite{friedrich2017}:
    \begin{equation}\label{eq:DecayRateGeneralForm}
        A_\mathrm{i}
        =
        \sum_\mathrm{f}
        A_{\mathrm{i}\to\mathrm{f}}
    \,, \
        A_{\mathrm{i}\to\mathrm{f}}
        =
        \frac{\EH}{\hbar} \frac{4}{3} \alpha^2 \bar{E}^3_\mathrm{i,f}
        \sum_{u=x,y,z}
        \abs*{
            \langle\bm{\Psi}_\mathrm{f},r_u\slash\aB\bm{\Psi}_\mathrm{i}\rangle
        }^2
    \,.
    \end{equation}
    In terms of oscillator strengths one obtains
    \begin{equation}\label{eq:SpontaneousRate}
        A_\mathrm{i}
        =
        \frac{2\alpha^3\EH}{\hbar}
        \sum_{n',l',j'}
        \bar{E}^2_\mathrm{i,f} (-f_{\mathrm{i}\to\mathrm{f}})
    \,.
    \end{equation}
    The rate $A_\mathrm{i}$ in Eq.~\eqref{eq:SpontaneousRate} is the spontaneous
    emission rate, which is only non-zero for transitions where a photon is
    emitted and the final state's energy is lower than the energy of the initial
    state, i.e., $E_\mathrm{i}>E_\mathrm{f}$. Note that this restricts the
    quantum numbers in the sum.

    At sufficiently high temperatures the spectral density of the environment
    cannot be neglected. To approximate the spectral density of the environment
    one can employ Plank's famous formula \cite{plank1901} for the spectral
    density of a black-body. This yields an average number of photons at
    frequency $\omega$ and temperature $T$:
    \begin{equation}
        n(\omega,T)
        =
        \frac{1}{\exp(\hbar\omega\slash(\kBT))-1}
    \,.
    \end{equation}
    Here, $k_\mathrm{B}$ is Boltzmann's constant. Compared to the low
    temperature regime, the additional photons at temperature $T$ stimulate the
    emission of photons and thereby the decay of the initial state to
    lower-energy final states. Even depopulation of the initial state is
    possible by the absorption of photons from the environment. This process
    drives transitions to final states which are higher in energy than the
    initial state. According to Einstein's rate equations \cite{einstein1917}
    the total rate for depopulation of the initial state is
    \begin{equation}\label{eq:LifetimeFormula}
        \frac{\hbar}{\EH}\Gamma
        =
        2\alpha^3
        \sum_{n',l',j'}
        \bar{E}^2_\mathrm{i,f} (-f_{\mathrm{i}\to\mathrm{f}})
        (1+n(E_\mathrm{i,f}\slash\hbar,T))
        +
        2\alpha^3
        \sum_{n',l',j'}
        \bar{E}^2_\mathrm{f,i} f_{\mathrm{i}\to\mathrm{f}}
        n(E_\mathrm{f,i}\slash\hbar,T)
    \,,
    \end{equation}
    where each sum is restricted to quantum numbers for which each summand is
    positive. This entails that the first sum runs only over final states with
    lower energy compared to the initial state; this describes the emission of
    photons. The second sum runs solely over final states with larger energy
    than the initial state; this corresponds to the case of photon absorption.

    \clearpage

    %%% SECTION
    \section{Manual}
    \label{sec:Manual}

    Reference manual for the \alkcalc library functions. Prototypes for the
    following functions are defined in the header `interface/alkcalc.h'.

    There is always a decision to be made when dealing with half-integer quantum
    numbers numerically. Here we agree on the following: say $q$ is a quantum
    number that may take half-integer values, e.g., $j$, $m_j$. To `clean up'
    the numerical value we do $2q\mapsto[2.0\times{q}+0.5]$. The `floor'
    function is defined as $[x]=k\in\mathbb{Z}$ ($x\in\mathbb{R}$), where $k$ is
    the unique integer such that $x\in[k,k+1)$. In practice, this means that for
    any $q\in[-0.75,-0.25)$ we use $q=-1\slash{2}$, for $q\in[-0.25,0.25)$ we
    use $q=0$, for $q\in[0.25,0.75)$ we use $q=1\slash{2}$, and so on, in both
    the negative and positive direction. In the following manual, this
    convention is applied to all quantum numbers that can take half-integer
    values, unless stated otherwise.

    \sffamily%
    \pagestyle{plain}%
    \vspace*{12mm}\par%

    \begin{sdocfun}
    {
        alkcalc\tus{Enlsj}
    }{
        Eigenenergies in units of Hartree. The requested eigenenergy is read
        from the corresponding data file in the directory `data' and returned.
    }
        \sdocarg{char *species}{
            String to specify atom/ion species, e.g., \texttt{1H} for Hydrogen,
            or \texttt{88SR+} for the $\mathrm{^{88}Sr^+}$ ion.
        }
        \sdocarg{int32\tus{t} n}{
            Principal quantum number.
        }
        \sdocarg{int32\tus{t} l}{
            Orbital angular momentum quantum number $l=0,1,\dots,n-1$.
        }
        \sdocarg{double s}{
            Electron spin quantum number. Not an argument, as $s=1\slash{2}$ is
            always assumed.
        }
        \sdocarg{double j}{
            Total angular momentum quantum number
            $j=\abs*{l-1\slash{2}},l+1\slash{2}$.
        }
        \sdocret{double E}{
            $E$ is the eigenenergy specified by the quantum numbers $n$, $l$,
            $s=1\slash{2}$, and $j$.
        }
    \end{sdocfun}

    \begin{sdocfun}
    {
        alkcalc\tus{fnlsj}
    }{
        Numerical equivalent to $f_{n,l,s,j}$ in Eq.~\eqref{eq:fRconversion}.
        The data is read from the file containing the requested state. This data
        is stored at a user specified location. The location where \alkcalc will
        search for the data can be set in the file `interface/alkcalc.h'. The
        result is owned by the caller and should be freed with
        \texttt{alkcalc\tus{state}\tus{free}}.
    }
        \sdocarg{char result}{
            Character for deciding if discretisation data is returned; see
            RETURN down below. If \texttt{result} is \texttt{f} (full), the full
            result is returned, and if \texttt{result} is \texttt{p} (partial),
            the discretisation data is not returned.
        }
        \sdocarg{char *species}{
            String to specify atom/ion species, e.g., \texttt{1H} for Hydrogen,
            or \texttt{88SR+} for the $\mathrm{^{88}Sr^+}$ ion.
        }
        \sdocarg{int32\tus{t} n}{
            Principal quantum number.
        }
        \sdocarg{int32\tus{t} l}{
            Orbital angular momentum quantum number $l=0,1,\dots,n-1$.
        }
        \sdocarg{double s}{
            Electron spin quantum number. Not an argument, as $s=1\slash{2}$ is
            always assumed.
        }
        \sdocarg{double j}{
            Total angular momentum quantum number
            $j=\abs*{l-1\slash{2}},l+1\slash{2}$.
        }
        \sdocret{alkcalc\tus{state} *f}{
            \texttt{f} contains eight members: \texttt{N}, \texttt{dim},
            \texttt{t}, \texttt{h}, and \texttt{fnlsj}.\par%
            \texttt{int32\tus{t} N}: Number of points, $N$, to discretise the
            interval $[0,t_\mathrm{max}]$ (see Sec.~\ref{sec:Theory}).\par%
            \texttt{int32\tus{t} n}: Principal quantum number.\par%
            \texttt{int32\tus{t} l}: Orbital angular momentum quantum
            number.\par%
            \texttt{double j}: Total angular momentum quantum number.\par%
            \texttt{int32\tus{t} dim}: Dimension of eigenproblem, i.e.,
            $\text{\texttt{dim}}=N-2$ [see
            Eq.~\eqref{eq:expansionOfState}].\par%
            \texttt{double *t}: Array of length $N$, containing the mesh
            $(t_k)_{k=0,\dots,N-1}$ (see Sec.~\ref{sec:Theory}).\par%
            \texttt{double *h}: Array of length $N-1$, containing the steps
            $(h_k)_{k=1,\dots,N-1}$ (see Sec.~\ref{sec:Theory}).\par%
            \texttt{double *fnlsj}: \texttt{dim}-dimensional eigenvector
            $\bm{f}$ [see Eq.~\eqref{eq:expansionOfState}].
        }
    \end{sdocfun}

    \begin{sdocfun}
    {
        alkcalc\tus{rp}
    }{
        Radial matrix element of $r^p$, $p\in\mathbb{R}$ [see
        Eq.~\eqref{eq:radialMatrixElement}]. For $\abs*{p}\gtrsim10$ there is
        the risk of numerical over or underflow of the 64-bit floating-point
        arithmetic. It is the user's responsibility to ensure that the integral
        associated to the desired matrix element exists. In case it does not,
        the returned value is wrong and no error is thrown. $\mathrm{^1H}$
        (Hydrogen atom) and $\mathrm{^2He^+}$ (Helium ion) can be used to
        benchmark atoms and ions, respectively.
    }
        \sdocarg{char *species}{
            String to specify atom/ion species, e.g., \texttt{1H} for Hydrogen,
            or \texttt{88SR+} for the $\mathrm{^{88}Sr^+}$ ion.
        }
        \sdocarg{int32\tus{t} nb}{
            Principal quantum number of bra.
        }
        \sdocarg{int32\tus{t} lb}{
            Orbital angular momentum quantum number $l=0,1,\dots,n-1$ of bra.
        }
        \sdocarg{double sb}{
            Electron spin quantum number. Not an argument, as $s=1\slash{2}$ is
            always assumed.
        }
        \sdocarg{double jb}{
            Total angular momentum quantum number of bra, $j$.
        }
        \sdocarg{double p}{
            Power of radius operator; cf. Eq.~\eqref{eq:radialMatrixElement}.
            For instance, $p=0$ yields the scalar product of the states. For
            $p=1$ and $p=2$, respectively, dipole and quadrupole matrix elements
            are computed. Generally, $p\in\mathbb{R}$.
        }
        \sdocarg{int32\tus{t} nk}{
            Principal quantum number of ket.
        }
        \sdocarg{int32\tus{t} lk}{
            Orbital angular momentum quantum number $l=0,1,\dots,n-1$ of ket.
        }
        \sdocarg{double sk}{
            Electron spin quantum number. Not an argument, as $s=1\slash{2}$ is
            always assumed.
        }
        \sdocarg{double jk}{
            Total angular momentum quantum number of ket, $j$.
        }
        \sdocret{double rp}{
            $r^{n,l,s,j;n',l',s',j'}_p$ in Eq.~\eqref{eq:radialMatrixElement} in
            units of $\aB^p$, where $\aB$ is Bohr's radius. Quantum numbers with
            a prime correspond to the bra, and the ones without one, to the ket.
        }
    \end{sdocfun}

    \begin{sdocfun}
    {
        alkcalc\tus{cj1m1j2m2jmj}
    }{
        Clebsch-Gordan coefficients, $C^{j_1,m_1,j_2,m_2}_{j,m_j}=\braket*{j_1,
        m_1;j_2,m_2}{j_1,j_2,j,m_j}$, where $j$ is the total angular momentum
        composed of the two individual ones $j_1$ and $j_2$. To fix the phase
        uniquely, we adopt the Condon-Shortley sign convention, that is, the
        Clebsch-Gordan coefficients with $m_j=j_1+j_2$ are positive. In
        particular, for the case $j_1=l$ and $j_2=s$,
        \begin{equation}\label{eq:definitionCGC}
            C^{l,m_l,s,m_s}_{j,m_j}
            =
            \langle
            Y_{l,m_l}\otimes\bm{\chi}_{s,m_s},
            \bm{\Phi}_{l,s,j,m_j}
            \rangle
        \,,
        \end{equation}
        with the states $Y_{l,m_l}\otimes\bm{\chi}_{s,m_s}$ of the uncoupled
        basis and the coupled basis states $\bm{\Phi}_{l,s,j,m}$; see
        Eqs.~(\ref{eq:sphericalHarmonics},\,\ref{eq:twoSpinsors}), and
        Eq.~\eqref{eq:coupledBasis}, respectively.  The return type (see RETURN
        down below) allows to extract the requested coefficient symbolically as
        well as numerically. The function returns zero in case the arguments do
        not make sense, e.g., if $j_1<m_1$; or if $j_1$ is half-integer and
        $m_1$ is integer; or if $j$ does not lay between (including bounds)
        $\abs*{j_1-j_2}$ and $\abs*{j_1+j_2}$.
    }
        \sdocarg{double j1}{
            Angular momentum quantum number of first factor.
        }
        \sdocarg{double m1}{
            Magnetic quantum number of first factor.
        }
        \sdocarg{double j2}{
            Angular momentum quantum number of second factor.
        }
        \sdocarg{double m2}{
            Magnetic quantum number of second factor.
        }
        \sdocarg{double j}{
            Total angular momentum quantum number.
        }
        \sdocarg{double mj}{
            Magnetic quantum number of total angular momentum.
        }
        \sdocret{alkcalc\tus{cg} c}{
            \texttt{c} contains three members: \texttt{sign},
            \texttt{numerator}, and \texttt{denominator}.\par%
            \texttt{int8\tus{t} sign}: Sign, $\sigma$, of the Clebsch-Gordan
            coefficient.\par%
            \texttt{int64\tus{t} numerator}: Numerator, $\nu$, of the
            Clebsch-Gordan coefficient.\par%
            \texttt{int64\tus{t} denominator}: Denominator, $\mu$, of the
            Clebsch-Gordan coefficient.\par%
            \noindent
            The requested Clebsch-Gordan coefficient is given by
            $\sigma\sqrt{\nu\slash\mu}$. Euklid's algrithm is employed to reduce
            $\nu$ and $\mu$. This representation allows to obtain Clebsch-Gordan
            coefficients symbolically as well as numerically.
        }
    \end{sdocfun}

    \begin{sdocfun}
    {
        alkcalc\tus{YlmlXsms}
    }{
        Angular factor of the full state in the uncoupled basis, i.e., the
        two-component spinor $Y_{l,m_l}(\vartheta,\varphi)\bm{\chi}_{s,ms}$.
        The input quantum numbers are checked for inconsistencies and an error
        is thrown in case one is detected. The angles $\vartheta$ and $\varphi$
        have the ranges $[0,\pii]$ and $[0,2\pii]$, respectively. This condition
        is checked up to floating point machine precision (type \texttt{double})
        and an error will be thrown in case either condition is not obeyed.
    }
        \sdocarg{int32\tus{t} l}{
            Orbital angular momentum quantum number $l=0,1,\dots,n-1$.
        }
        \sdocarg{int32\tus{t} ml}{
            Magnetic quantum number of orbital angular momentum.
        }
        \sdocarg{double s}{
            Electron spin quantum number. Not an argument, as $s=1\slash{2}$ is
            always assumed.
        }
        \sdocarg{double ms}{
            Magnetic quantum number of electron spin, $m_s$. For any $m_s<0$
            the value $m_s=-1\slash{2}$ is used internally, and for $m_s\geq0$
            the value $m_s=1\slash{2}$.
        }
        \sdocarg{double theta}{
            Polar angle (Zenitwinkel), $\vartheta$, in range $[0,\pii]$.
        }
        \sdocarg{double phi}{
            Azimuthal angle (Azimut), $\varphi$, in range $[0,2\pii]$.
        }
        \sdocret{alkcalc\tus{spinor} s}{
            \texttt{s} has two components: \texttt{u} and \texttt{v}.\par%
            \texttt{u} is the spin-up ($m_s=1\slash{2}$) component.\par%
            \texttt{d} is the spin-down ($m_s=-1\slash{2}$) component.\par%
        }
    \end{sdocfun}

    \begin{sdocfun}
    {
        alkcalc\tus{Philsjmj}
    }{
        Angular factor of the full eigenstate in the coupled (fine-structure)
        basis, i.e., the two-component spinor
        $\bm{\Phi}_{l,s,j,mj}(\vartheta,\varphi)$ [see
        Eq.~\eqref{eq:coupledBasis}]. The input quantum numbers are checked for
        inconsistencies and an error is thrown in case one is detected. The
        angles $\vartheta$ and $\varphi$ have the ranges $[0,\pii]$ and
        $[0,2\pii]$, respectively. This condition is checked up to floating
        point machine precision (type \texttt{double}) and an error will be
        thrown in case either condition is not obeyed. To compute the resulting
        two-component spinor, Clebsch-Gordon coefficients are needed, for which
        the sign convention detailed in Eq.~\eqref{eq:definitionCGC} is
        employed.
    }
        \sdocarg{int32\tus{t} l}{
            Orbital angular momentum quantum number $l=0,1,\dots,n-1$.
        }
        \sdocarg{double s}{
            Electron spin quantum number. Not an argument, as $s=1\slash{2}$ is
            always assumed.
        }
        \sdocarg{double j}{
            Total angular momentum quantum number
            $j=\abs*{l-1\slash{2}},l+1\slash{2}$.
        }
        \sdocarg{double mj}{
            Magnetic quantum number of total angular momentum, $m_j$.
        }
        \sdocarg{double theta}{
            Polar angle (Zenitwinkel), $\vartheta$, in range $[0,\pii]$.
        }
        \sdocarg{double phi}{
            Azimuthal angle (Azimut), $\varphi$, in range $[0,2\pii]$.
        }
        \sdocret{alkcalc\tus{spinor} s}{
            \texttt{s} has two components: \texttt{u} and \texttt{v}.\par%
            \texttt{u} is the spin-up ($m_s=1\slash{2}$) component.\par%
            \texttt{d} is the spin-down ($m_s=-1\slash{2}$) component.\par%
        }
    \end{sdocfun}

    \begin{sdocfun}
    {
        alkcalc\tus{fitof}
    }{
        Oscillator strength for the transition between an initial state
        $\bm{\Psi}_\mathrm{i}$ and a final state $\bm{\Psi}_\mathrm{f}$, as
        defined in Subsec.~\ref{subsec:OscillatorStrengths}, that is, it was
        already summed over the magnetic quantum number of the final state.
    }
        \sdocarg{char *species}{
            String to specify atom/ion species, e.g., \texttt{1H} for Hydrogen,
            or \texttt{88SR+} for the $\mathrm{^{88}Sr^+}$ ion.
        }
        \sdocarg{int32\tus{t} ni}{
            Principal quantum number of the initial state
            $\bm{\Psi}_\mathrm{i}$.
        }
        \sdocarg{int32\tus{t} li}{
            Orbital angular momentum quantum number $l=0,1,\dots,n-1$ of the
            initial state $\bm{\Psi}_\mathrm{i}$.
        }
        \sdocarg{double si}{
            Electron spin quantum number. Not an argument, as $s=1\slash{2}$ is
            always assumed.
        }
        \sdocarg{double ji}{
            Total angular momentum quantum number of the initial state
            $\bm{\Psi}_\mathrm{i}$, $j$.
        }
        \sdocarg{int32\tus{t} nf}{
            Principal quantum number of the final state $\bm{\Psi}_\mathrm{f}$.
        }
        \sdocarg{int32\tus{t} lf}{
            Orbital angular momentum quantum number $l=0,1,\dots,n-1$ of the
            final state $\bm{\Psi}_\mathrm{f}$.
        }
        \sdocarg{double sf}{
            Electron spin quantum number of the final state
            $\bm{\Psi}_\mathrm{f}$. Not an argument here, since $s=1\slash{2}$
            is always assumed.
        }
        \sdocarg{double jf}{
            Total angular momentum quantum number of the final state
            $\bm{\Psi}_\mathrm{f}$, $j$.
        }
        \sdocret{double fitof}{
            $f_\mathrm{i\to{f}}$ is the oscillator strength of the selected
            electronic transition, summed over $m'_j$ of the final state. Note
            that because of this sum, it is the same for all $m_j$ of the
            initial state.
        }
    \end{sdocfun}

    \begin{sdocfun}
    {
        alkcalc\tus{tau}
    }{
        Lifetime of the fine-structure state $\bm{\Psi}_{n,l,s,j,m_j}$ in
        nanoseconds. Note that, since the oscillator-strengths [see
        Eq.~\eqref{eq:OscillatorStrengthSimplified}] do not depend on $m_j$,
        also the lifetimes do not depend on it [see
        Eq.~\eqref{eq:LifetimeFormula}].
    }
        \sdocarg{double T}{
            Temperature of black-body excitation spectrum in Kelvin.
        }
        \sdocarg{char *species}{
            String to specify atom/ion species, e.g., \texttt{1H} for Hydrogen,
            or \texttt{88SR+} for the $\mathrm{^{88}Sr^+}$ ion.
        }
        \sdocarg{int32\tus{t} n}{
            Principal quantum number.
        }
        \sdocarg{int32\tus{t} dn}{
            Range for principal quantum number. At zero temperature, this
            argument is not used. For finite temperature, the absorption
            processes in Eq.~\eqref{eq:LifetimeFormula} allow `decay' to states
            with higher energy. In principle, one needs to take into account all
            transitions to these states, of which there are infinitely many.
            Here, we only take into account \texttt{dn} levels above
            $\bm{\Psi}_{n,l,s,j,m_j}$ to truncate the series, for both
            $l'=l\pm1$. In other words, the sum in
            Eq.~\eqref{eq:LifetimeFormula}, taking into account absorption,
            contains $2\times\text{\texttt{dn}}$ terms for $l>0$ and \texttt{dn}
            for $l=0$.
        }
        \sdocarg{int32\tus{t} l}{
            Orbital angular momentum quantum number $l=0,1,\dots,n-1$.
        }
        \sdocarg{double s}{
            Electron spin quantum number. Not an argument, as $s=1\slash{2}$ is
            always assumed.
        }
        \sdocarg{double j}{
            Total angular momentum quantum number
            $j=\abs*{l-1\slash{2}},l+1\slash{2}$.
        }
        \sdocret{double tau}{
            $\tau$ is the lifetime of the fine-structure state
            $\bm{\Psi}_{n,l,s,j,m_j}$ (independent of $m_j$) in nanoseconds.
        }
    \end{sdocfun}

    %%% BIBLIOGRAPHY
    \newpage
    \renewcommand\bibname{References}%
    \printbibliography

\end{document}
